% Minimal retained preamble for lane-p2 cards taken from
% sources/logic-fl02-conservation-of-ramsey-s-theorem-f/pinned-original-source.tex
% (Conservation of Ramsey's theorem for pairs, arXiv:2402.11616v3).
%
% Copied verbatim from the author's own file:
%   line 2      : \documentclass{article}
%   lines 16-25 : the author's theorem declarations
%   lines 46-88 : the author's macro definitions for the reverse-mathematics principles
%   line 91     : \card
%   lines 104-107: the natural sums and products, which rely on mathabx (source line 103)
%   line 116    : \qt
%
% Omitted author preamble lines and why (nothing was renamed or redefined):
%   line 3 (graphicx), 7 (url), 9 (cancel): no figure, url or cancel use in the retained spans;
%   line 6 (cleveref): the assembler adds its own cleveref machinery, a second load clashes;
%   line 8 (\usepackage[utf8x]{inputenc}): inputenc with the utf8x option fights the forced
%       xelatex engine;
%   lines 28-37 (the rep@theorem machinery and \newreptheorem{maintheorem}): the retained spans
%       use no rep-theorem environment;
%   lines 40-43 (xcolor and the three author comment macros): not used by the retained spans;
%   lines 93-99 (the LS2 font encoding block for \smashtimes based on stix): the host TeX Live has
%       no stix directory and the author himself never uses the result -- line 105 maps \ntimes to
%       \dottimes, and line 110 shows \smashtimes already commented out in the source.
% mathabx is retained because \dotplus, \dottimes and \dotdiv come from it.
\documentclass{article}
\usepackage{amsmath,amssymb,amsthm}
\usepackage{mathabx}

\newtheorem{theorem}{Theorem}
\numberwithin{theorem}{section}
\newtheorem{maintheorem}[theorem]{Main Theorem}
\newtheorem{lemma}[theorem]{Lemma}
\newtheorem{question}[theorem]{Question}
\newtheorem{proposition}[theorem]{Proposition}
\newtheorem{statement}[theorem]{Statement}
\newtheorem{remark}[theorem]{Remark}
\newtheorem{definition}[theorem]{Definition}
\newtheorem{corollary}[theorem]{Corollary}

\newcommand{\RCA}{\mathsf{RCA}}
\newcommand{\WKL}{\mathsf{WKL}}
\newcommand{\SRT}{\mathsf{SRT}}
\newcommand{\EM}{\mathsf{EM}}
\newcommand{\SEM}{\mathsf{SEM}}
\newcommand{\ADS}{\mathsf{ADS}}
\newcommand{\SCAC}{\mathsf{SCAC}}
\newcommand{\SADS}{\mathsf{SADS}}
\newcommand{\COH}{\mathsf{COH}}
\newcommand{\ACA}{\mathsf{ACA}}
\newcommand{\qvdash}{\operatorname{{?}{\vdash}}}
\newcommand{\nqvdash}{\operatorname{{?}{\nvdash}}}
\newcommand{\BSig}{\mathsf{B}\Sigma^0}
\newcommand{\ISig}{\mathsf{I}\Sigma^0}
\newcommand{\LSig}{\mathsf{L}\Sigma^0}
\newcommand{\BPi}{\mathsf{B}\Pi^0}
\newcommand{\IPi}{\mathsf{I}\Pi^0}
\newcommand{\LPi}{\mathsf{L}\Pi^0}
\newcommand{\IDel}{\mathsf{I}\Delta^0}
\newcommand{\BCDe}{\mathsf{BC}\Delta^0}
\newcommand{\WF}{\mathsf{WF}}
\newcommand{\RT}{\mathsf{RT}}
\newcommand{\BME}{\mathsf{BME}}
\newcommand{\TT}{\mathsf{TT}}
\renewcommand{\O}{\mathcal{O}}
\renewcommand{\L}{\mathcal{L}}
\newcommand{\C}{\mathcal{C}}
\newcommand{\R}{\mathcal{R}}
\renewcommand{\S}{\mathcal{S}}
\newcommand{\T}{\mathcal{T}}
\renewcommand{\P}{\mathcal{P}}
\newcommand{\M}{\mathcal{M}}
\newcommand{\Nc}{\mathcal{N}}
\renewcommand{\P}{\mathcal{P}}
\newcommand{\Q}{\mathcal{Q}}
\newcommand{\NN}{\mathbb{N}}
\newcommand{\PP}{\mathbb{P}}
\newcommand{\uh}[0]{{\upharpoonright}}
\newcommand{\card}{\operatorname{card}}
\newcommand{\nplus}{\dotplus}
\newcommand{\ntimes}{\dottimes}
\newcommand{\nminus}{\dotdiv}
\newcommand{\nsum}{\dot{\sum}}
\def\qt#1{``#1''}
\title{Well-foundedness of an ordinal carries over to its natural product with any finite coefficient}
\author{Quentin Le Houérou, Ludovic Levy Patey, Keita Yokoyama}
\date{arxiv:2402.11616v3; CC BY 4.0}
\begin{document}
\maketitle
\noindent\textit{Editorial note.} The counted unit is the article's own Lemma lem:wf-product at source lines 411-413 together with the proof environment that begins on the line after the statement ends (415-423); the pairing is the ordinary adjacent one, and the proof contains the author's own commented-out alternative argument at source lines 418-422, which is retained exactly as the author left it.\par\medskip
\noindent\textit{Editorial note.} The retained context at source line 390 is the author's own announcement of the two combinatorial lemmas and of the natural (Hessenberg) operations; the retained context at source lines 392-398 is his own definition of the natural sum and of alpha ntimes k; the retained context at source line 400 is his own comparison of the natural product with the usual ordinal product, which the proof uses.\par\medskip
\noindent\textit{Editorial note.} The retained preamble reproduces verbatim the author's document class at source line 2, his theorem declarations at source lines 16-25, his macro definitions at source lines 46-88 and 91, his natural-operation macros at source lines 104-107 and his quotation macro at source line 116. The LS2 font-encoding block at source lines 93-99, which declares \textbackslash{}smashtimes through the stix fonts, is not retained because the host TeX Live has no stix directory and because the author never uses that symbol: source line 105 defines \textbackslash{}ntimes as \textbackslash{}dottimes and source line 110 shows the alternative already commented out. mathabx (source line 103) is retained because \textbackslash{}dotplus and \textbackslash{}dottimes come from it. The cleveref load at source line 6 and the inputenc load with the utf8x option at source line 8 are not retained, for the assembler's own cleveref machinery and the forced xelatex engine respectively. Nothing was renamed, redefined or added.\par\medskip
\noindent\textit{Editorial note.} No retained line contains a citation or a cross-reference, so the material carries no bibliography and needs no navigation pointer.\par\medskip
\noindent\textit{Editorial note.} The article's conservation theorems, its lemma lem:shore-blocking-wf, its lemmas lem:coh-wf-forcing-jump and lem:coh-wf-forcing-wf, and its proposition prop:coh-wf are not counted here and are not used as premises.\par\medskip
\noindent\textit{Allowed standard prerequisites (editorial).}\begin{itemize}
\item Natural sum and natural product of ordinals below epsilon\_0. For ordinals alpha and beta below epsilon\_0 written with a common family of exponents, the natural sum is obtained by adding the coefficients of like powers of omega, and alpha ntimes k is the natural sum of alpha with itself k times; the article's own definition is retained at source lines 392-398, together with the comparison alpha ntimes k at most alpha times omega at source line 400.. The counted lemma concludes WF(alpha ntimes k) from WF(alpha), so both the natural product and its comparison with the usual product must be available.
\item Theories RCA\_0 and the well-foundedness statement WF. RCA\_0 is the standard base theory of reverse mathematics, and WF(alpha) states that the ordinal alpha is well founded; a model M satisfies WF(alpha) when it has no M-infinite descending sequence of ordinals below alpha; the article uses these notations throughout, for instance when it defines the Pi\textasciicircum{}1\_1 proof-theoretic ordinal in its introduction.. The counted lemma is stated for a model M of RCA\_0 plus WF(alpha) and concludes WF(alpha ntimes k) inside the same model.
\end{itemize}\medskip
Before working on $\Pi^1_1$-conservation of~$\WKL$ and $\COH$, we prove two combinatorial lemma which say that Shore blocking comes for free for preserving $\WF(\alpha)$ for any ordinal $\alpha \leq \epsilon_0$. Instead of using the standard sums and products over ordinals, the natural (Hessenberg) sums and products will be more convenient:

\begin{definition}[Natural sum and product]
Let $\alpha$ and $\beta$ be two ordinals less than $\epsilon_0$. Let $\alpha = \omega^{\gamma_1}n_1 + \dots + \omega^{\gamma_k}n_k$ and $\beta = \omega^{\gamma_1}m_1 + \dots + \omega^{\gamma_k}m_k$ (We allow the $n_i$ and $m_i$ to be equal to $0$ in order to write $\alpha$ and $\beta$ using the same exponents $\gamma_i$). We define the natural sum of $\alpha$ and $\beta$ (denoted by $\alpha \nplus \beta$) to be the following ordinal (It does not depend on the choice of the family $(\gamma_i)$) : 

$$\omega^{\gamma_1}(n_1 + m_1) + \dots + \omega^{\gamma_k}(n_k + m_k)$$

We also define $\alpha \ntimes k$ to be equal to be the natural sum of $\alpha$ with itself $k$ times. %and $\alpha \ntimes \omega = \omega^{\gamma_1+1}n_1 + \dots + \omega^{\gamma_k+1}n_k$ (whereas $\alpha \times \omega = \omega^{\gamma_1 + 1}n_1$)
\end{definition}

Note that $\alpha \ntimes k \leq \alpha \times \omega$, where $\alpha \times \omega = \sup_{k < \omega} \alpha \times k = \omega^{\gamma_1 + 1}n_1$ is the usual ordinal product.

\begin{lemma}\label[lemma]{lem:wf-product}
Let~$\M = (M,S) \models \RCA_0 + \WF(\alpha)$ for $\alpha$ an ordinal less than $\epsilon_0$. Then $\M \models \WF(\alpha \ntimes k)$ for every $k \in M$.
\end{lemma}

\begin{proof}
Since $\alpha \ntimes k \leq \alpha \times \omega$ for every $k \in M$ and $\RCA_0$ proves that the product of two well-orders is a well-order, hence $\RCA_0 \vdash \WF(\alpha) \rightarrow \WF(\alpha \times \omega)$, we get the desired result.

% Consider an M-infinite decreasing sequence of ordinals $(\alpha_s)_{s \in M}$ all less than $\alpha \ntimes k$. Each $\alpha_s$ can be written as $\alpha_{s,0} \# \dots \# \alpha_{s,k-1}$ where we choose $\alpha_{s,0}$ to be the biggest ordinal less than or equal to $\alpha$ satisfying $\alpha_{s,0} \leq \alpha_s$, we then choose $\alpha_{s,1}$ to be the biggest ordinal less than or equal to $\alpha$  satisfying $\alpha_{s,0} \# \alpha_{s,1} \leq \alpha_s$...

% Since $(\alpha_s)$ is decreasing, the sequences $(\alpha_{s,i})$ are non-decreasing for every $i < k$, and for every $s$, there is at least one $i < k$ such that $\alpha_{s,i} > \alpha_{s+1,i}$.

% But, by $\RT^1_k$, which follows from $\BSig_2$, there exists some~$i < k$ and an $M$-infinite set~$H \in \M$ such that $\alpha_{s,i} > \alpha_{s+1,i}$ for all $s \in H$. Then $\{\alpha_{s,i} : s \in H \}$ is an $M$-infinite decreasing sequence in~$\alpha$, contradicting our assumption.
\end{proof}

\end{document}
